\documentclass[12pt]{article}
\usepackage[utf8]{inputenc}
\usepackage[T1]{fontenc}
\usepackage{hyperref}
\title{Sadri, Oraon Language Resources for K-12 Curriculum Development}
\author{Compiled by Claude Sonnet 5.5 (claude-sonnet-5-5) for Andr\'as Kornai}
\date{}
\begin{document}
\maketitle

\section{Introduction}
Oraon Sadri (ISO 639-3 sdr) is the Indo-Aryan variety of Sadri used by Oraon communities, notably in Bangladesh. It is distinct from the Dravidian Kurukh, the traditional Oraon language. Sadri is a regional lingua franca in eastern India. It is also called Nagpuri in its literary form, which has its own ISO code, sck. Almost nothing was found for the sdr variety specifically. The resources below mostly concern Sadri and Nagpuri in general and may not match Bangladeshi Oraon usage.

\section{Primary Educational Resources}
No primers or textbooks for Oraon Sadri in Bangladesh were found, and no mother-tongue education programme for it was confirmed.

\section{Linguistic and Reference Materials}
\begin{itemize}
\item Peterson and Baraik, a grammar of Chotanagpuri Sadri, the Indo-Aryan lingua franca of eastern central India (a catalogue listing of the book). It describes Sadri in India, not Bangladesh. \url{https://www.exoticindiaart.com/book/details/grammar-of-chotanagpuri-sadri-indo-aryan-lingua-franca-of-eastern-central-india-hbh466}
\item Webonary site for Sadri. Its content and variety were not checked. \url{https://www.webonary.work/sadri/?p=9444}
\item Wikipedia article on the Nagpuri language. \url{https://en.wikipedia.org/wiki/Nagpuri_language}
\item Simple English Wikipedia article on Sadri. \url{https://simple.wikipedia.com/wiki/Sadri_language}
\end{itemize}

\section{Religious and Community Resources}
\begin{itemize}
\item Scripture Earth page for sck (Nagpuri/Sadri), listing New Testament and audio material. \url{https://www.scriptureearth.org/00eng.php?iso=sck}
\item Banglapedia entry on the Oraon of Bangladesh. \url{https://en.banglapedia.org/index.php/Oraon,_The}
\item University of Dhaka paper, ``The Oraons of Bangladesh: An Anthropological'' study, which discusses Sadri use in the community. \url{https://sanskrit.du.ac.bd/wp-content/uploads/2025/07/5.-The-Oraons-of-Bangladesh-An-Anthropological-converted.pdf}
\item Journal article record on the Oraon (Bangladesh Economic Association site). \url{https://bea-bd.org/site/article-details/180}
\item Joshua Project page on Kurux Oraon in Bangladesh. \url{https://www.joshuaproject.net/index.php/people_groups/14210/BG}
\item Wikipedia article on the Kurukh people. \url{https://en.wikipedia.org/wiki/Kurukh_people}
\end{itemize}

\section{Digital and Audio Resources}
Scripture Earth (above) lists Gospel audio and film for Sadri (sck). No Bangladesh-specific audio or software was found.

\section{Curriculum Development Context}
Sources report that Sadri has declined in family and other domains among Bangladeshi Oraons, and that some have moved to Bengali. Nagpuri is commonly written in Devanagari, and Bangladeshi Oraon practice may differ. A curriculum would need an orthography decision (Bengali or Devanagari script), and would need to decide whether to adapt Indian Sadri or Nagpuri materials. Community fieldwork in Bangladesh would also be needed.
\end{document}